\documentclass[IB]{TFUOC}%IB: CASTELLANO, CAT: CATALÁN, ENG: ANGLÈS
\graphicspath{{figuras/}{plantilla-uoc/}{sistema-diseno/}}

\usepackage[bibliografia]{tfm-estilo}
\addbibresource{referencias.bib}
\usepackage{multicol}
\usepackage{booktabs}
\usepackage{colortbl}
\usepackage{makecell}
\usepackage{xcolor}
\usepackage{tabularx}
\usepackage{float}
\usepackage{pgfgantt}
\usepackage[spanish]{babel}
\usepackage{hyperref}

%Introducción de datos del trabajo
\title{Sistema automatizado de clasificación de la accesibilidad de las vías urbanas a partir de imágenes aéreas de alta resolución}
\titcrt{Clasificación de accesibilidad urbana}
\author{Víctor Pariente González}
\date{24 de septiembre de 2026}


\nomPDC{Sergi Sanz Bravo}
\nomPRA{Antonio Ruiz Falcó Rojas}
\titulac{Máster en Ciencia de Datos}
\area{Área 3}
\idioma{Castellano}
\credits{12}
\parcla{accesibilidad, segmentación de imagen, extracción geométrica, planificación urbana}

\licenc{ccBy}
%Posibles licencias
%ccByNcNd
%ccByNcSa
%ccByNc
%ccByNd
%ccBySa
%ccBy
%GNU
%copyright


%%%%%%%%%%
% Resumen en el idioma

\abstractidioma{
El diseño de las ciudades no es accesible para todas las personas. En España, en 2020, había un 9,5\% de personas con algún tipo de discapacidad reconocida, que sumado a las personas mayores o con movilidad reducida temporal representan un aproximado 30\% de la población.

La legislación estatal obliga a las ciudades a modificar su entramado urbano para que sea accesible para toda la población, mediante la ejecución de planes de accesibilidad universal, que persiguen el diseño para todos teniendo en cuenta la diversidad funcional.

Sin embargo, el proceso para planificar y diseñar los planes es manual, lento y costoso, debido a que obliga a los equipos técnicos a desplazarse in situ para realizar mediciones.

Mediante una metodología de investigación basada en diseño, se porpone un sistema automatizado de clasificación de la accesibilidad de las vías urbanas a partir de imágenes aéreas de alta resolución, mediante segmentación semántica de imagen y extracción geométrica de medidas.

El objetivo es utilizar los resultados de las predicciones del sistema para facilitar a los ayuntamientos el diagnóstico y elaboración de los planes de actuación, así como identificar las zonas más prioritarias sobre las que actuar.}

% Resumen en inglés.
\abstractenglish{
City design is not accessible to everyone. In Spain, in 2020, 9.5\% of the population had a recognized disability. When combined with older adults and people with temporary reduced mobility, this represents approximately 30\% of the population.

National legislation requires cities to adapt their urban streets to ensure accessibility for every person through universal accessibility plans, which promote design for all while accounting for functional diversity.

However, the planning and drafting process for these plans is manual, slow, and costly, as technical teams are required to conduct on-site field measurements.

Using a Design Research methodology, this work proposes an automated system for classifying urban street accessibility from high aerial resolution images using semantic image segmentation and geometric measurement extraction.

The goal is to leverage the system's predictive output to help local councils diagnose accessibility issues, draft action plans, and prioritize the areas for intervention.}

\begin{document}

\estructura

\tableofcontents

\listoffigures

\listoftables




\chapter{Introducción}
\section{Contexto y justificación del trabajo}
Las ciudades tienen una \textbf{deuda histórica con la diversidad humana} de sus habitantes. Hasta décadas recientes, el entorno urbano se ha \textbf{diseñado para la ciudadanía <<normativa>>}, definía como una persona adulta joven, ágil y sin ningún tipo de limitación física, sensorial o cognitiva. La irrupción del \textbf{vehículo privado} en el día a día de las personas añadió un nuevo elemento predominante en las vías urbanas (calzadas, aparcamientos). Esto eliminó espacios para los viandantes, facilitando entornos con \textbf{aceras estrechas, desniveles o problemas de continuidad} que excluyen a varios sectores de la sociedad en el ejercicio de sus derechos y disfrute de la ciudad.
\\

Cualquier persona puede pasar a lo largo de su vida de ser <<normativa>> a tener necesidades específicas de movilidad. Por esa razón, la \textbf{accesibilidad universal} no es una demanda exclusiva de la ciudadanía con discapacidad permanente, sino una condición que \textbf{afecta a toda la sociedad} en distintas etapas de la vida.
\\

En España, según información de 2020, \textbf{4,38 millones de habitantes tenían algún tipo de diversidad funcional o discapacidad} (física, visual, auditiva, cognitiva o relacional), lo que representa aproximadamente un 9,5\% de la población. Este porcentaje se va acentuando con la edad, alcanzando valores superiores al \textbf{30\% en las personas mayores de 65 años} (colectivo que representa un 20\% del total y en constante aumento por el envejecimiento demográfico). No se pueden olvidar tampoco las personas con \textbf{movilidad reducida temporal} (lesiones, uso de muletas o silla de ruedas) ni las \textbf{familias con carritos de bebé}.
\\

Todos estos colectivos representan más de un 30\% de la ciudadanía, y reflejan que la problemática de la accesibilidad no solo impacta en ciertos colectivos, sino en amplios grupos de personas en el momento actual, y que \textbf{nos afectará en mayor o menor medida a todos}.
Citando a \textit{PREDIF}:
\begin{quote}
«La accesibilidad es imprescindible para el 10\% de la población, necesaria para el 40\% y cómoda para el 100».
\end{quote}

De esta necesidad, y con el objetivo de tratar de tirar abajo la deuda accesible de las ciudades, se han redactado varias leyes en la legislación estatal española, de las comunidades autónomas y municipal.
El marco estatal de referencia es la \textbf{Orden TMA/851/2021 de 23 de julio}, que legisla sobre las \textbf{condiciones básicas de accesibilidad} y no discriminación, y los derechos de utilización de los espacios públicos urbanos.
\\

En esta normativa se impone como mínimo a las administraciones municipales:
\begin{itemize}
    \item La obligación de elaborar los \textbf{Planes de Accesibilidad Universal (PAU)}, planes estratégicos para identificar el estado de las calles, programar y priorizar las intervenciones y comprobar su cumplimiento.
    \item Impulsar el concepto de \textbf{accesibilidad universal}, no solo pensando en la población con problemas de movilidad, sino atendiendo a todos los tipos de diversidad.
    \item Establecer los \textbf{requisitos mínimos} para considerar las vías urbanas como accesibles, bajo el concepto de \textbf{Itinerario Peatonal Accesible (IPA)}.
\end{itemize}


A pesar de esta legislación, según los diagnósticos realizados por el IMSERSO, el estado actual de la accesibilidad de las ciudades es bastante negligente. En el año 2005, más del \textbf{16\% de aceras eran totalmente inaccesibles} por tener anchos inferiores a lo establecido por ley, y peor aún, más del \textbf{63\% de los itinerarios peatonales tenían obstáculos físicos} que dificultan la movilidad de los viandantes. Según el informe de Índice de Accesibilidad e Inclusión de la Fundación ONCE (2022), el \textbf{37\%} de las necesidades de accesibilidad en la vía urbana \textbf{siguen incumpliendo la ley}.
\\

La ejecución de los planes de accesibilidad sufre de unos procedimientos excesivamente lentos y costosos, debido a que las inspecciones del equipo técnico deben hacerse de manera presencial y con materiales específicos, lo que conlleva un:
\begin{itemize}
    \item \textbf{Coste económico elevado:} por la mano de obra, los desplazamientos y el uso de material, especialmente en localidades más pequeñas.
    \item \textbf{Rápida caducidad de los diagnósticos:} al ser tan lentos, no se adaptan al crecimiento dinámico de las ciudades, como obras o cambios en el mobiliario urbano.
    \item \textbf{Procesamiento lento de los datos:} una vez obtenidas las mediciones, los datos se tienen que subir a las plataformas web, en algunos casos inexistentes.
\end{itemize}

La necesidad de acelerar los planes de accesibilidad es el punto de partida de este proyecto. Para ello, se plantea la \textbf{utilización de imágenes aéreas de alta resolución} como fuente de datos para el diagnóstico de las vías urbanas, utilizando técnicas de inteligencia artificial, construyendo un \textbf{sistema que procese de manera automática} la evaluación de las calles de cada ciudad.
\\

Este planteamiento permitirá \textbf{reducir la carga de trabajo, la lentitud y los costes} que conlleva la inspección \textit{in situ} del procedimiento actual, \textbf{impactando} positivamente en el conjunto de la \textbf{ciudadanía}.
\\

\section{Motivación personal}
Mi interés personal a la hora de seleccionar este trabajo parte principalmente de dos motivos. En primer lugar, buscaba realizar un proyecto con\textbf{ carga social} y que tuviera \textbf{aplicación real y humana}, intentando impactar en la mejora de la vida de las personas.
\\

En segundo lugar, quería que estuviese \textbf{relacionado con la accesibilidad}, debido a que dentro de mi \textbf{círculo más cercano} cuento con \textbf{personas con diversidad funcional}, como la hipoacusia o problemas de movilidad, y gracias a su activismo y su relación con plataformas de accesibilidad como \textit{PREDIF} o \textit{Music For All}, me ha permitido conocer más de cerca el trabajo que realizan estas \textbf{asociaciones}, y el largo camino que queda por recorrer, además de conocer gente con diferentes diversidades.
\\

Una vez solicitado el proyecto, tuve la mala suerte de \textbf{sufrir una lesión física} en el pie, que me ha obligado durante semanas a \textbf{utilizar silla de ruedas y muletas para desplazarme}.
\\

Esta vivencia personal me ha permitido \textbf{comprender mucho mejor todos los obstáculos} con los que las ciudades nos impiden desplazarnos con autonomía e independencia.
Además, he podido comprobar que la \textbf{accesibilidad} no es algo estático, sino \textbf{dinámico y dependiente del tiempo}, lo que hace que calles en principio accesibles se transformen en problemáticas.
\\

Un ejemplo es la presencia de cubos de basura depositados en la acera por la noche en el centro de Madrid, que reducen de manera drástica el ancho de las aceras, la disposición de mobiliario urbano, como las marquesinas de autobús de manera que se impida el paso incluso para los viandantes; o el \textbf{incivismo de la sociedad}.
\\

Esta experiencia personal me permitirá ayudar a identificar características y variables a la hora de ejecutar el proyecto.
\\

\section{Objetivos del trabajo}

\subsection{Objetivo general}

El objetivo general del trabajo es el \textbf{diseño e implementación de un sistema automatizado de clasificación de parámetros de accesibilidad del viario urbano}, a partir de  imágenes aéreas de alta resolución y basado en técnicas de exportación geométrica, aprendizaje profundo y visión artificial. El sistema estará integrado en un \textbf{visor web interactivo y exportable}, sirviendo como herramienta de soporte a la realización de Planes de Accesibilidad Universal (PAU), el fomento de la participación ciudadana y la transparencia de las administraciones locales.

A continuación se detalla una \textbf{lista de objetivos principales}, que comprenden el \textbf{producto mínimo viable} (MVP, siglas de Minimum Valuable Product) del trabajo; y una lista de objetivos secundarios de valor añadido.

\subsection{Objetivos principales (MVP)}

Este listado comprende los objetivos planificados para cubrir obligatoriamente, siendo la columna vertebral del proyecto.
\\

\textbf{OP1. Identificar necesidades y requisitos de accesibilidad universal}

Analizar e identificar las necesidades de movilidad y accesibilidad en el espacio público de los entornos urbanos para toda la ciudadanía —teniendo en cuenta la diversidad funcional (física, auditiva, visual, cognitiva, relacional), así como al resto de la población— con el fin de definir los requisitos que garanticen la accesibilidad universal en el entorno urbano, y cumplan el principio de diseño para toda la población.
\\

\textbf{OP2. Identificar y definir las métricas e indicadores de evaluación de la accesibilidad}

Analizar la normativa vigente (estatal, autonómica y local) para identificar el conjunto de indicadores cuantitativos y cualitativos (tales como la anchura libre de paso, o criterios necesarios para plataformas únicas) y establecer las fórmulas matemáticas de evaluación de la accesibilidad de las vías urbanas.
\\

\textbf{OP3. Entrenar modelos de \textit{deep learning} para la segmentación semántica de imágenes}

Entrenar un modelo con técnicas de aprendizaje profundo (\textit{deep learning}) y visión por computador sobre un conjunto de imágenes aéreas de alta resolución, para la identificación y clasificación semántica automatizada de aceras, calzadas, aparcamientos y obstáculos en las ciudades.
\\

\textbf{OP4. Desarrollar un módulo de extracción geométrica de métricas de accesibilidad}

Implementar un procedimiento de geometría computacional en \textit{Python} que, a partir de los segmentos obtenidos en el punto anterior, calcule las magnitudes necesarias para realizar los cálculos de indicadores y métricas de accesibilidad, como la anchura de aceras o intersecciones.
\\

\textbf{OP5. Predecir y validar el algoritmo con un conjunto de datos}

Ejecutar el algoritmo sobre un conjunto de imágenes de una ciudad piloto, y evaluar de manera cuantitativa la precisión de las predicciones obtenidas por el modelo, cruzando los datos contra los registros previamente calculados sobre el terreno, empleando las métricas de evaluación según el estado del arte.
\\

\textbf{OP6. Desarrollar un algoritmo de alineación con el callejero}
\\

Integrar los datos obtenidos con los datos cartográficos del callejero de las ciudades a predecir.
\\

\textbf{OP7. Construir un visor GIS interactivo con exportación de datos}

Diseñar e implementar una aplicación web que permita visualizar mediante un mapa la accesibilidad de las calles por tramos, así como la descarga de los datos en formatos abiertos (CSV o JSON).
\\

\subsection{Objetivos secundarios}

Representan las metas adicionales o de extensión al trabajo principal del proyecto. Debido a las limitaciones temporales, y al tiempo de dedicación disponible, la consecución de estos objetivos es un valor añadido.
\\

\textbf{OS1. Construir un portal web informativo y de participación ciudadana}

Diseñar un portal web web con orientación pública para divulgar los planes y políticas de accesibilidad universal, habilitando un canal interactivo de consulta y participación ciudadana.
\\

\textbf{OS2. Incluir un catálogo de cumplimiento de accesibilidad por municipio}

Integrar en el portal ciudadano un catálogo de accesibilidad municipal que muestre información detallada de los planes de accesibilidad en marcha, así como el visor interactivo de cada localidad.
\\

\textbf{OS3. Habilitar la integración con estándares cartográficos de geoportales}

Estructurar los datos obtenidos tras la predicción como servicios web de mapas, permitiendo una integración directa con Sistemas de Información Geográfica (SIG) o geoportales.
\\

\section{Impacto en sostenibilidad, ético-social y de diversidad}
\label{s:etic}

Una vez analizados los Objetivos de Desarrollo Sostenible (ODS) de la Agenda 2030 de la ONU, se enumeran a continuación los impactos positivos, negativos y su mitigación, en relación al presente trabajo fin de máster.

\subsection{Sostenibilidad y huella ecológica}

Los ODS a los que más se contribuye en el proyecto son: \textbf{ODS 11 (\textit{Sustainable Cities and Communities})} y \textbf{ODS 9 (\textit{Industry, Innovation and Infraestructure})}.
\\

En cuanto a los \textbf{impactos positivos}, se establecen:
\begin{itemize}
    \item \textbf{Sustitución de trabajos in situ}: el mecanismo tradicional de auditoría de accesibilidad incluye un desplazamiento continuo de vehículos, para que el personal técnico se desplace a cada calle, así como material técnico para poder realizar las mediciones. La utilización del sistema automático usando imágenes aéreas reduce de manera significativa los desplazamientos, permitiendo una reducción de las emisiones indirectas CO2.
    \item \textbf{Optimización de los recursos públicos}: además de la reducción de costes de gasolina por los desplazamientos, la posibilidad de ordenar los tramos de calles menos accesibles permite a los ayuntamientos la priorización de obras de reparación en los tramos más críticos, optimizando las intervenciones en los puntos más necesarios.
\end{itemize}

Sin embargo, como \textbf{impacto negativo} se encuentra:
\begin{itemize}
    \item \textbf{Consumo energético por entrenamiento de los modelos de inteligencia artificial}: la utilización de algoritmos de aprendizaje automático va asociada una huella de carbono digital.
\end{itemize}

La \textbf{mitigación} de este punto se basa en la preferencia del uso de modelos ya preentrenados, utilizando técnicas de \textit{transfer learning}, lo que hace reducir el tiempo de uso eléctrico y computacional a la hora de entrenar el modelo.

\subsection{Comportamiento ético y responsabilidad social}

El ODS que aplica a este punto es: \textbf{ODS 16 (Peace, Justice and Strong Institutions)}
\\

Como \textbf{impactos positivos}, principalmente en gobernanza abierta y no discriminación, se incluyen:
\begin{itemize}
    \item \textbf{Democratización del dato mediante estándares abiertos}: los resultados de las predicciones del sistema se publican en formatos abiertos (como CSV), por lo que se elimina cualquier dependencia de aplicaciones propietarias de pago, o la necesidad de contratar consultoría tecnológica privada. Al incluir también un portal para que municipios pequeños puedan utilizar la herramienta, se permite la igualdad de condiciones.
    \item \textbf{Transparencia institucional y auditoría social}: la implantación de un catálogo abierto de índice de accesibilidad por municipio permite a la ciudadanía y a los organizmos a visibilizar el grado de cumplimiento con la ley de cada localidad, habilitando una auditoría pública.
\end{itemize}

\subsection{Diversidad, género y derechos humanos}

Uno de los objetivos principales del trabajo es el enfoque accesible universal, por lo que los ODS a los que se vincula son \textbf{ODS 10 (Reduced Inequalities)} y \textbf{ODS 5 (Gender Equality)}.
\\

Sobre los \textbf{impactos positivos}, se encuentran:
\begin{itemize}
    \item \textbf{Inclusión y diversidad funcional global}: los beneficios de los planes de accesibilidad no están únicamente pensados para colectivos particulares, sino que están pensados para todo el conjunto diverso de la población. Por lo que no se benefician unicamente en personas con discapacidad permanente, también disfrutan de los planes las personas mayores, personas con discapacidad temporal, familias o el resto de la población.
    \item \textbf{Derecho a la autonomía y disfrute de la ciudad}: el proyecto facilita la implantación de actuaciones de eliminación de barreras en la vía pública y tiene una visión global de la diversidad, por lo que permite una mayor autonomía personal, así como un mejor derecho de disfrute de la ciudad. 
    \item \textbf{Perspectiva de género en la movilidad urbana}: el aumento de autonomía en la diversidad funcional implica una menor necesidad de ayuda para los desplazamientos, algo que impacta positivamente en la habitual mayor carga de cuidados de las mujeres.
\end{itemize}

\section{Enfoque y método seguido}

Utilizando la referencia \textit{Researching Information Systems and Computing} de \textit{Oates}, y analizando la tipología del trabajo, se ha elegido como \textbf{estrategia principal} la de \textbf{investigación basada en diseño}, debido a que se va a realizar la entrega de un sistema informático con modelos de aprendizaje automático para la evaluación de accesibilidad, usando una estrategia adicional de estudio de caso para comparar el sistema con el visor de accesibilidad del ayuntamiento de Albacete.
\\

\subsection{Estrategia principal: investigación basada en diseño (\textit{Design Research})}

Las fases que comprende esta estrategia son:
\begin{enumerate}
    \item \textbf{Identificación del problema (\textit{Awareness})}: identificación del problema grave de las ciudades que no cumplen con criterios de accesibilidad única, análisis de la legislación y buenas prácticas recomendadas por organizaciones accesibles, análisis de los planes de actuación y comprobación de las limitaciones actuales de los procedimientos manuales de verificación de parámetros de accesibilidad, con visitas in situ del personal técnico de los ayuntamientos.
    \item \textbf{Diseño (\textit{Suggestion}):} análisis de literatura científica con el estado del arte del diseño de soluciones para el análisis de imágenes aéreas urbanas y extracción de medidas geométricas. Realización de una propuesta de diseño tecnológica con técnicas de aprendizaje profundo y visión por computador, así como la interfaz web de usuario de visualización con mapas y exportación.
    \item \textbf{Desarrollo (\textit{Development})}: implementación del sistema informático en Python, incluyendo los modelos y los artefactos web necesarios para cumplir con el diseño planificado.
    \item \textbf{Evaluación (\textit{Evaluation})}: validación cuantitativa de la precisión del modelo, comparándolo con los datos calculados manualmente, lo que se conoce como \textit{Ground Truth} en la ciencia de datos.
    \item \textbf{Conclusión (\textit{Conclusion})}: documentación de los resultados obtenidos, de los aprendizajes obtenidos tras la realización del trabajo, así como enumeración de futuras líneas de ampliación.
\end{enumerate}

\subsection{Estrategia adicional: caso de estudio (\textit{Case Study})}

Tomando como base el \textbf{visor GIS del ayuntamiento de Albacete}, se puede realizar una evaluación del comportamiento del sistema, utilizando imágenes aéreas de la localidad de Albacete, y \textbf{comparando los resultados} obtenidos con el portal web de la localidad.

De esta manera, se podrá demostrar la calidad del sistema, y servirá para demostrar el uso real de la herramienta en los ayuntamientos.

\subsection{Metodología de ingeniería de software}
Se plantea la ejecución del diseño e implementación de código como un proyecto de ingeniería de software, con tecnologías de codigo abierto, herramientas de inteligencia artificial, y de gestión de proyectos.

\begin{enumerate}
    \item Metodología
        \begin{itemize}
            \item \textbf{Desarrollo ágil con metodología Kanban:} organización del trabajo utilizando tableros Kanban en GitHub projects.
            \item \textbf{Principios de código limpio y control de versiones:} implementación del código siguiendo los principios de \textit{Clean Code}, utilizando control de versiones con el flujo de \textit{Git Flow}.
        \end{itemize}
    \item Stack tecnológico
        \begin{itemize}
            \item \textbf{Lenguaje de programación:} Python, el lenguaje estándar para los trabajos de ciencia de datos.
            \item \textbf{Librerías de \textit{deep learning} y visión por computador:} uso de librerías de Python como PyTorch o Torchvision.
            \item \textbf{Librerías de tratamiento espacial de datos:} utilización de librerías como GeoPandas, NumPy, exportación en formatos abiertos (GeoJson/CSV), y tecnologías estándar de desarrollo web (HTML5, CSS, Javascript).
        \end{itemize}
    \item Recursos hardware
        \begin{itemize}
            \item \textbf{Entorno con GPU:} se necesita utilizar un entorno con aceleración gráfica por hardware (GPU), ya sea en entorno local o remoto, para poder entrenar los modelos.
        \end{itemize}
    \item Herramientas
        \begin{itemize}
            \item \textbf{Procesador de textos:} se utilizará Overleaf como plataforma de procesador de textos LaTeX para la redacción de la memoria.
            \item \textbf{Repositorio de código:} GitHub como repositorio de código fuente y control de versiones.
            \item \textbf{Seguimiento ágil del proyecto:} se utilizará GitHub Projects/Trello para controlar el seguimiento del trabajo, utilizando épicas e historias de usuario.
        \end{itemize}
    \item Plataformas de formación
        \begin{itemize}
            \item Se utilizarán \textbf{plataformas de aprendizaje online} como DataCamp o Coursera para consolidar y refrescar los conceptos de Python y sus librerías de ciencia de datos con impacto en el proyecto.
        \end{itemize}    
    \item Utilización de la inteligencia artificial  
\end{enumerate}


\section{Planificación del trabajo}

Se ha realizado una \textbf{planificación en base a la carga de 12 créditos ECTS}, equivalentes a 300 horas de dedicación, distribuyendo la carga de trabajo en las 15 semanas lectivas disponibles.

La estructuración del trabajo se ha realizado utilizando \textbf{paquetes de trabajo (WP)}, que se han basado principalmente en las entregas establecidas en las pruebas de integración continua de la UOC, incluyendo una serie de \textbf{hitos} para reflejar la evolución incremental del trabajo tanto en el estudio del estado del arte y en el diseño e implementación del mismo.

Además, se han incluido paquetes de trabajo para los \textbf{objetivos secundarios}, planificados inicialmente como opcionales.

\subsection{Calendario del trabajo}

En primer lugar, se muestra en la siguiente tabla la división de las tareas a realizar durante el trabajo, descompuestas en \textbf{6 paquetes de trabajo principales}, indicando el número de créditos ETCS y horas, y a qué entrega de evaluación continua se corresponden. Para el hito de implementación (WP3), se ha realizado una \textbf{división inicial en subfases}, que podría verse actualizada una vez se consolide la investigación del estado del arte (WP2).

\begin{table}[H] %
  \centering
  \caption{Planificación del proyecto y distribución de carga ECTS}
  \label{tab:planificacion_tfm}
  \scriptsize
  \setlength{\extrarowheight}{1pt}

  \begin{tabularx}{\textwidth}{l X c c c}
    \toprule
    \textbf{WP} & \textbf{Tarea y descripción detallada} & \textbf{Carga} & \textbf{Fechas} & \textbf{UOC} \\
    \midrule
    WP1 & \textbf{Definición y objetivos:} planteamiento inicial, problemática, alineación de objetivos, metodología y planificación temporal. & \makecell{0.6 ECTS \\ (15 h)} & 24 sep - 04 oct & PEC1 \\
    \midrule
    WP2.1 & \textbf{Análisis legal:} estudio de legislación y normativa sobre accesibilidad e infraestructura urbana, extrayendo requisitos. & \makecell{0.6 ECTS \\ (15 h)} & 05 oct - 12 oct & PEC2 \\
    WP2.2 & \textbf{Estado del arte:} revisión de literatura sobre aprendizaje profundo para segmentación semántica y extracción geométrica. & \makecell{1.0 ECTS \\ (25 h)} & 13 oct - 25 oct & PEC2 \\
    \midrule
    WP3.1 & \textbf{Preparación de datos:} recopilación, limpieza y preprocesamiento de imágenes aéreas y etiquetado de la red viaria. & \makecell{0.8 ECTS \\ (20 h)} & 19 oct - 03 nov & PEC3 \\
    WP3.2 & \textbf{Entrenamiento IA:} entrenamiento y ajuste fino de redes neuronales (transfer learning) para segmentar aceras y calzadas. & \makecell{1.2 ECTS \\ (30 h)} & 25 oct - 24 nov & PEC3 \\
    WP3.3 & \textbf{Extracción geométrica:} algoritmos de visión por computador para derivar métricas físicas reales desde las máscaras. & \makecell{1.0 ECTS \\ (25 h)} & 03 nov - 24 nov & PEC3 \\
    WP3.4 & \textbf{Clasificación:} motor lógico de clasificación del nivel de accesibilidad cruzando medidas con requisitos legales. & \makecell{0.8 ECTS \\ (20 h)} & 11 nov - 24 nov & PEC3 \\
    WP3.5 & \textbf{Validación e iteración:} evaluación frente al conjunto de prueba, análisis de errores y ajuste continuo de hiperparámetros. & \makecell{0.6 ECTS \\ (15 h)} & 13 nov - 05 dic & PEC3 \\
    WP3.6 & \textbf{Integración con el callejero:} realización de la integración del modelo con los datos geográficos del callejero de una ciudad, para poder incorporar los resultados del análisis al visor. & \makecell{0.6 ECTS \\ (15 h)} & 18 nov - 03 dic & PEC3 \\
    WP3.7 & \textbf{Visor GIS e interfaz:} desarrollo de un visor geográfico interactivo para la representación gráfica sobre mapas de la red urbana. & \makecell{0.6 ECTS \\ (15 h)} & 20 nov - 05 dic & PEC3 \\
    \rowcolor{gray!15}
    --- & \textbf{Descanso:} pausa planificada sin carga lectiva ni entregas asignadas. & \makecell{0.0 ECTS \\ (0 h)} & 02 dic - 07 dic & --- \\
    WP3.8 & \textbf{Portal web ciudadano:} portal web orientado al ciudadano integrando un canal para consultas e incidencias urbanas. & \makecell{0.6 ECTS \\ (15 h)} & 08 dic - 13 dic & PEC3 \\
    WP3.9 & \textbf{Catálogo municipio:} observatorio municipal web para la consulta y análisis comparativo de índices de accesibilidad. & \makecell{0.6 ECTS \\ (15 h)} & 08 dic - 14 dic & PEC3 \\
    WP3.10 & \textbf{Exportación GIS:} publicación de resultados en estándares de servicios web de mapas (WMS/WFS) para interoperabilidad. & \makecell{0.4 ECTS \\ (10 h)} & 11 dic - 15 dic & PEC3 \\
    \midrule
    WP4 & \textbf{Redacción memoria:} consolidación y redacción de la memoria final del TFM, estructurando metodología y resultados. & \makecell{1.6 ECTS \\ (40 h)} & 16 dic - 27 dic & PEC4 \\
    \midrule
    WP5 & \textbf{Presentación vídeo:} elaboración del guion, soporte visual y grabación del recurso audiovisual requerido para la defensa. & \makecell{1.0 ECTS \\ (25 h)} & 28 dic - 05 ene & PEC5 \\
    \midrule
    WP6 & \textbf{Defensa oral:} ensayo, preparación de respuestas a posibles objeciones del tribunal y exposición oral final del proyecto. & \makecell{1.0 ECTS \\ (25 h)} & 06 ene - 21 ene & PEC6 \\
    \bottomrule
  \end{tabularx}
\end{table}

El cronograma del proyecto, incluyendo tanto los paquetes de trabajo como hilos de entrega, se detalla en el siguiente diagrama de Gantt. En color azul se encuentran las entregas principales, mientras que en color verde se indican las entregas secundarias.

\begin{figure}[H]
  \centering
  \scriptsize

  \definecolor{mainbar}{RGB}{51,102,204}
  \definecolor{secbar}{RGB}{46,125,50}
  \definecolor{breakbar}{RGB}{180,180,180}
  \definecolor{milestonecolor}{RGB}{211,47,47}

  \begin{ganttchart}[
      vgrid={*6{dotted}, *1{dashed}},
      hgrid,
      x unit=1.05mm, % Ancho temporal optimizado para A4
      time slot format=isodate,
      title height=1,
      y unit chart=0.42cm,
      y unit title=0.5cm,
      bar/.append style={fill=mainbar, draw=none},
      bar height=0.5,
      milestone/.append style={fill=milestonecolor, draw=none},
      milestone height=0.6,
      bar label font=\tiny\sffamily, % Fuente limpia para los nombres
      milestone label font=\tiny\itshape\color{milestonecolor},
      title label font=\tiny\bfseries,
      canvas/.append style={draw=black!30}
    ]{2026-09-24}{2027-01-21}

    % Escala de tiempo
    \gantttitlecalendar{year, month=shortname} \\

    % WP1
    \ganttbar{WP1 Definición y objetivos}{2026-09-24}{2026-10-04} \\
    \ganttmilestone{Entrega PEC1}{2026-10-04} \\

    % WP2
    \ganttbar{WP2.1 Análisis legal}{2026-10-05}{2026-10-12} \\
    \ganttbar{WP2.2 Estado del arte}{2026-10-13}{2026-10-25} \\
    \ganttmilestone{Entrega PEC2}{2026-10-25} \\

    % WP3
    \ganttbar{WP3.1 Preparación de datos}{2026-10-19}{2026-11-03} \\
    \ganttbar{WP3.2 Entrenamiento IA}{2026-10-25}{2026-11-24} \\
    \ganttbar{WP3.3 Extracción geométrica}{2026-11-03}{2026-11-24} \\
    \ganttbar{WP3.4 Clasificación}{2026-11-11}{2026-11-24} \\
    \ganttbar{WP3.5 Validación e iteración}{2026-11-13}{2026-12-05} \\
    \ganttbar{WP3.6 Integración con el callejero}{2026-11-18}{2026-12-03} \\
    \ganttbar{WP3.7 Visor GIS e interfaz}{2026-11-20}{2026-12-05} \\

    % Descanso
    \ganttbar[bar/.append style={fill=breakbar}]{Descanso}{2026-12-02}{2026-12-07} \\

    % Secundarios
    \ganttbar[bar/.append style={fill=secbar}]{WP3.8 Portal web ciudadano}{2026-12-08}{2026-12-13} \\
    \ganttbar[bar/.append style={fill=secbar}]{WP3.9 Catálogo municipio}{2026-12-08}{2026-12-14} \\
    \ganttbar[bar/.append style={fill=secbar}]{WP3.10 Exportación GIS}{2026-12-11}{2026-12-15} \\
    \ganttmilestone{Entrega PEC3}{2026-12-15} \\

    % WP4
    \ganttbar{WP4 Redacción memoria}{2026-12-16}{2026-12-27} \\
    \ganttmilestone{Entrega PEC4}{2026-12-27} \\

    % WP5 & WP6
    \ganttbar{WP5 Presentación vídeo}{2026-12-28}{2027-01-05} \\
    \ganttmilestone{Entrega PEC5}{2027-01-05} \\
    \ganttbar{WP6 Defensa oral}{2027-01-06}{2027-01-21} \\
    \ganttmilestone{Entrega PEC6 / Final}{2027-01-21}

  \end{ganttchart}
  \caption{Diagrama de Gantt del proyecto.}
  \label{fig:gantt_tfm}
\end{figure}


\subsection{Identificación de riesgos}
Con el fin de asegurar el cumplimiento de la planificación, así como la consecución de los hitos de entrega académicos, se ha realizado una matriz de evaluación de posibles riesgos que pueden tener impacto en la elaboración del trabajo, incluyendo qué probabilidad tienen de producirse, qué impacto pueden tener en el proyecto, y el plan de contingencia a seguir si se producen.

Los principales riesgos identificados han estado centrados en la incentidumbre de la calidad de entrenamiento y predicción de los modelos de aprendizaje profundo, así como la priorización de objetivos principales frente a los secundarios.

\begin{table}[H]
  \centering
  \caption{Matriz de identificación de riesgos y planes de contingencia}
  \label{tab:identificacion_riesgos}
  \scriptsize
  \setlength{\extrarowheight}{3pt}
  
  \begin{tabularx}{\textwidth}{>{\raggedright\arraybackslash}p{2.8cm} c >{\raggedright\arraybackslash}X >{\raggedright\arraybackslash}X}
    \toprule
    \textbf{Factor de riesgo} & \textbf{Prob. / Impacto} & \textbf{Impacto en el proyecto} & \textbf{Plan de contingencia} \\
    \midrule
    
    \textbf{Dificultad en el modelo de segmentación semántica} & 
    Media/Alto & 
    Constituye uno de los principales riesgos del proyecto al ser el punto de inicio del resto del sistema. Problemas posibles: menor precisión de los modelos en tramos que estén cubiertos por vegetación, toldos, o aglomeraciones, así como dificultades a la hora de hiperparametrizar los modelos. & 
    Filtrado inicial de imágenes priorizando las que tienen mejor visualización, segmentación específica de tramos con vegetación y aplicación de técnicas de \textit{data augmentation} para generar imágenes sin vegetación. \\
    \midrule
    
    \textbf{Tiempos elevados de entrenamiento de modelos} & 
    Media/Medio & 
    Los trabajos de ciencia de datos se basan en la experimentación y prueba continua, por lo que el tiempo de entrenamiento constituye una variable importante en la consecución de objetivos. Con tiempos de entrenamiento elevados, se reduce el número de iteraciones que se pueden realizar y validar. & 
    Priorización del uso de \textit{transfer learning} sobre redes preentrenadas basándose en el estado del arte analizado y empleo de aceleración por GPU. \\
    \midrule
    
    \textbf{No llegar a tiempo para completar los objetivos secundarios} & 
    Media/Bajo & 
    En caso de que la planificación de los objetivos principales no se cumpla a tiempo, hay un riesgo de que se extienda el tiempo dedicado a estos objetivos, influyendo en el tiempo disponible para realizar los objetivos secundarios del proyecto. & 
    Dar prioridad a los hitos principales, extendiendo la dedicación en el cronograma y reevaluando cada semana el estado de los hitos secundarios en la planificación. \\
    \midrule
    
    \textbf{Indisponibilidad personal por vacaciones planificadas} & 
    Alta/Medio & 
    Se ha incluido en el calendario un período de descanso de una semana (del 2 al 7 de diciembre) por un período vacacional planificado, lo que puede afectar al tiempo disponible para la realización del proyecto. & 
    adelanto de la fase implementación (WP3) respecto a lo indicado por el calendario de entregas de la UOC, intentando completar con margen previo las entregas previas (WP1 y WP2) \\
    
    \bottomrule
  \end{tabularx}
\end{table}


\section{Breve sumario de productos obtenidos}

Los productos que se obtendrán tras la ejecución del proyecto se detallan a continuación, distinguiendo los entregablos del MVP de los secundarios:

\textbf{Productos del MVP}
\begin{enumerate}
    \item Conjunto de datos de imágenes aéreas de alta resolución.
    \item Repositorio en GitHub de código fuente del modelo de entrenamiento de la segmentación semántica de imágenes, así como una exportación del mismo.
    \item Repositorio y código fuente de los procedimientos de extracción geométrica, así como de los algoritmos de clasificación accesible.
    \item Repositorio y código fuente del portal web GIS interactivo.
    \item Archivos de datos con las predicciones obtenidas del modelo a partir del conjunto de imágenes iniciales.
    \item Memoria del trabajo.
\end{enumerate}

\textbf{Productos secundarios}
\begin{enumerate}
    \item Repositorio y código fuente del portal web ciudadano.
    \item Conjunto de datos del observatorio de transparencia municipal, así como el repositorio y código fuente del algoritmo de extracción de datos.
    \item Servicios web de mapas preparados para su consumo en geoportales.
\end{enumerate}

\section{Breve descripción de los otros capítulos de la memoria}

La estructura de la memoria tendrá los siguientes capítulos, coincidiendo en orden con las entregas estipuladas en el portal de la universidad:

\begin{enumerate}
    \item \textbf{Introducción y contexto del proyecto:} define el contexto y motivación del proyecto, una introducción a la normativa legal, el listado de objetivos, la metodología y la planificación del trabajo.
    \item \textbf{Características de vías accesibles. Marco legal y normativo:} define la legislación vigente en relación a las ciudades accesibles universales, así como las características de las vías accesibles en base a la ley, que servirán de base para los algoritmos.
    \item \textbf{Estado del arte:} contiene el análisis del estado el arte científico en cuanto a modelos de aprendizaje automático y algoritmos utilizados en problemas similares.
    \item \textbf{Arquitectura del sistema:} describe la arquitectura tecnológica del sistema, incluyendo detalles de cada uno de los módulos del mismo: algoritmos de \textit{deep learning}, de visión computacional, y clasificación de accesibilidad de calles.
    \item \textbf{Diseño e implementación de los algoritmos:} contiene las decisiones de diseño de los sistemas, las librerías y tecnológícas utilizadas, y los detalles y decisiones de implementación.
    \item \textbf{Resultados y evaluación:} detalla los resultados de las predicciones realizadas, lista las rutas a los sistemas y conjuntos de datos, y analiza la evaluación de la calidad del trabajo.
    \item \textbf{Conclusiones y trabajos futuros:} contiene las conclusiones finales del estudiante y el listado de líneas futuras de trabajos.
\end{enumerate}

\section{Declaración sobre el uso de inteligencia artificial generativa}


A continuación, se lista la utilización de inteligencia artificial generativa durante el desarrollo del trabajo:

\begin{itemize}
    \item \textbf{Facilitación de ingesta de legislación y literatura científica:} se ha usado NotebookLM para poder analizar y comprender mejor textos legales y científicos y generando artefactos como mapas mentales o infografías conceptuales, o para obtener las secciones más interesantes a leer de las fuentes seleccionadas.
    \begin{itemize}
        \item En concreto, se han incluido todos los textos legales del BOE que aplican al problema, empezando por la Orden TMA/851/2021 de 23 de julio, consultando qué otras legislaciones vigentes tienen relación con la citada ley, o si hay algúna legislación autonómica más restrictiva.
        \item Así como para realizar preguntas con algunas dudas de entendimiento de la ley, como por ejemplo: "¿Qué se considera como itinerario, un tramo de una calle que va de un tramo de calzada a otro? ¿O se cuenta la manzana?", o "¿Qué significa continuidad en un itinerario peatonal accesible, en que momento se puede considerar el tramo como no accesible?".
        \item También para facilitar la búsqueda de criterios legislativos que puedan utilizarse a partir de visión aérea de la ciudad, o resolver dudas sobre los mínimos de distancias y elementos permitidos en la vía pública.
    \end{itemize}
    \item \textbf{Recursos gráficos para la memoria:} con herramientas de inteligencia artificial generativa para incluir imágenes o esquemas que faciliten la lectura de la memoria y la enriquezcan.
    \begin{itemize}
        \item En primer lugar, se ha generado un sistema de diseño utilizando Claude Design, a partir de la plantilla de la memoria de la UOC, que incluye una paleta de colores, tipografías y componentes (gráficos, tablas, diagramas...) que permiten enriquecer visualmente la memoria y mantener un estilo coherente en todo el documento.
    \end{itemize}
    \item \textbf{Asistencia en formateo de textos en LaTeX:} como apoyo técnico para poder transformar textos escritos en textos formateados para LaTeX, o maquetación de tablas o diagramas.
    \item \textbf{Asistencia programación y refactorización de código:} se utilizarán asistentes de código como Claude o Copilot como ayuda en la escritura de \textit{scripts} en Python.
    \item \textbf{Diseño UI/UX y estilos del portal web:} con el fin de ofrecer una página web más atractiva para el usuario, se emplearán herramientas de inteligencia artificial que den soporte con el diseño y estilos de la web.
\end{itemize}

\chapter{Características de vías accesibles. Marco legal y normativo.}
\section{Marco legal y normativo}

\section{Características de vías accesibles}

\subsection{Itinerario Peatonal Accesible (IPA)}

\subsection{Pavimentos y elementos táctiles indicadores}

\subsection{Vados peatonales, pasos de peatones e isletas de refugio}

\subsection{Obras e intervenciones}

\subsection{Mobiliario urbano y terrazas}

\subsection{Calles de plataforma única}

\subsection{Señalización y elementos de información}

\chapter{Estado del arte}
\input{capitulos/03-estado-del-arte.tex}

\chapter{Arquitectura del sistema}
\input{capitulos/04-arquitectura-del-sistema.tex}

\chapter{Diseño e implementación de los algoritmos}
\input{capitulos/05-diseno-implementacion.tex}

\chapter{Resultados y evaluación}
\input{capitulos/06-resultados-evaluacion.tex}

\chapter{Conclusiones y trabajos futuros}
\section{Conclusiones}
Este capítulo tiene que incluir:
\begin{itemize}
\item Una descripción de las conclusiones del trabajo:
\begin{itemize}
    \item Una vez se han obtenido los resultados, ¿qué conclusiones se extraen?
    \item ¿Estos resultados son los esperados? ¿O han sido sorprendentes? ¿Por qué? 
\end{itemize}
\item Una reflexión crítica sobre el logro de los objetivos planteados inicialmente:
\begin{itemize}
    \item ¿Hemos logrado todos los objetivos? Si la respuesta es negativa, ¿por qué motivo?
\end{itemize}
\item Un análisis crítico del seguimiento de la planificación y metodología a lo largo del producto:
\begin{itemize}
    \item ¿Se ha seguido la planificación?
    \item ¿La metodología prevista ha sido suficientemente adecuada?
    \item ¿Ha habido que introducir cambios para garantizar el éxito del trabajo? ¿Por qué? 
\end{itemize}
\item De los impactos previstos en \ref{s:etic}, ético-sociales, de sostenibilidad y de diversidad, evaluar/mencionar si se han mitigado (si eran negativos) o si se han conseguido (si eran positivos). 
\item Si han aparecido impactos no previstos a \ref{s:etic}, evaluar/mencionar cómo se han mitigado (si eran negativos) o que han aportado (si eran positivos).
\item Las líneas de trabajo futuro que no se han podido explorar en este trabajo y han quedado pendientes.


%\item Una descripción de las conclusiones del trabajo: Qué lecciones se han aprendido del trabajo?
%\item Una reflexión crítica sobre el logro de los objetivos planteados inicialmente: Hemos conseguido todos los objetivos? Si la respuesta es negativa, por qué motivo?
%\item De los impactos previstos a la sección \ref{s:etic}, una evaluación, o al menos, mención, sobre si se han mitigado (si eran negativos) o si se han conseguido (si eran positivos). 
\end{itemize}

%\section{Líneas de futuro}
%Las líneas de trabajo futuro que no se han podido explorar en este trabajo y han quedado pendientes.

%\section{Seguimiento de la planificación}
%Un análisis crítico del seguimiento de la planificación y metodología a lo largo del trabajo: 
%\begin{itemize}
%    \item  Se ha seguido la planificación? 
%    \item La metodología prevista ha sido la adecuada? 
%    \item Ha habido que introducir cambios para garantizar el éxito del trabajo? ¿Por qué? 
%\end{itemize}

\section{Líneas futuras de trabajo}
Durante la fase de análisis de objetivos, se han tenido en cuenta las siguientes líneas de trabajo que pueden dar continuación y mejorar los hitos reflejados en el proyecto:
\\

\textbf{Línea 1. Construir un portal web informativo y de participación ciudadana}

Diseñar un portal web web con orientación pública para divulgar los planes y políticas de accesibilidad universal, habilitando un canal interactivo de consulta y participación ciudadana.
\\

\textbf{Línea 2. Incluir un catálogo de cumplimiento de accesibilidad por municipio}

Integrar en el portal ciudadano un catálogo de accesibilidad municipal que muestre información detallada de los planes de accesibilidad en marcha, así como el visor interactivo de cada localidad.
\\

\textbf{Línea 3. Habilitar la integración con estándares cartográficos de geoportales}

Estructurar los datos obtenidos tras la predicción como servicios web de mapas, permitiendo una integración directa con \gls{SIG} o geoportales.
\\

\textbf{Línea 4. Implementar un algoritmo de recomendación de plataformas únicas}

Diseñar un modelo prescriptivo que, a partir de las métricas extraídas, permita identificar vías inaccesibles que, según los criterios de la normativa legal vigente, sean candidatas para su transformación en plataformas únicas con prioridad peatonal.
\\

\textbf{Línea 5. Desarrollar un módulo abierto para la ingesta y procesamiento de imágenes}

Facilitar en el portal público un flujo para el proceso de imágenes aéreas o satelitales, habilitando a municipios pequeños y a la ciudadanía  la visualización y análisis de datos de sus localidades.
\\

\textbf{Línea 6. Incorporar en el modelo de accesibilidad factores temporales o dinámicos}

Existen factores temporales que, para acotar el alcance del proyecto no se han incorporado, relacionadas con el la característica dinámica de las ciudades. Entre otros factores encontrados se han encontrado: barreras y obstáculos en determinadas franjas horarias, como carga y descarga, aglomeración de personas en eventos, mercadillos, horario de depósito de cubos de basura en la calle, entrada y salida de colegios y trabajos, o la ocupación de aceras de los talleres mecánicos cuando están abiertos.
\\

\textbf{Línea 7. Mapa de rutas accesibles en las ciudades}

Modelar una red de caminos accesibles de un punto a otro de la ciudad, que permitan realizar un trayecto con un mínimo nivel de accesibilidad, permitiendo visualizar en un mapa la ruta, y la comparación con la ruta más corta que no cumple con los criterios, permitiendo medir la diferencia y sobreesfuerzo que se necesita para completar la misma ruta.
\\

\textbf{Línea 8. Sistema híbrido con imágenes a nivel de calle}

Realizar un sistema que contenga tanto modelos que analicen las imágenes aéreas, con imágenes a nivel de calle obtenidas a partir de softwares como Google Street View, o imágenes particulares, que permitan realizar incluir factores no visibles desde imágenes en dos dimensiones aéreas, como la altura de las aceras, la estructura de semáforos, o la calidad del pavimento.
\\

\chapter{Glosario}
\begin{multicols}{2}
\begin{description}
    \item[Accesibilidad universal:] condición de no discriminación de cualquier persona, permitiendo la autonomía propia, derechos y disfrute de su entorno, independientemente de tener diversidad funcional o no.
    \item[Ancho libre de paso:] banda de un itinerario peatonal que permite una circulación continua a las personas, con una dimensión mínima de 1,80 metros, sin ningún obstáculo en la totalidad de su recorrido (Orden TMA/851/2021, artículo 5.2.b).
    \item[Deep learning (aprendizaje profundo):] campo del machine learning basado en redes neuronales con múltiples capas, capaz de aprender a reconocer automáticamente elementos complejos en imágenes, optimizando la matriz de pesos del modelo.
    \item[Diseño para todos:] concepto de diseño por el que desde el principio se tiene en cuenta la diversidad de toda la población en materia de utilización y disfrute de los entornos.
    \item[Diversidad funcional:] concepto que define el conjunto de la sociedad con capacidades diversas (física, visual, auditiva, cognitiva o relacional) y que tiene en cuenta a todas las personas independientemente de sus capacidades. Sustituye al término discapacidad, con una connotación más social.
    \item[Extracción geométrica:] procedimiento que utiliza como entrada las máscaras de predicción de la segmentación semántica para obtener medidas reales (como metros) y calcular dimensiones a partir de operaciones vectoriales.
    \item[GitHub:] plataforma basada en Git que actúa como repositorio de código fuente, que se utiliza para el control de versiones y la gestión del ciclo de vida del proyecto.
    \item[GPU (Graphics Processing Unit):] unidad de procesamiento que permite la aceleración por hardware y el cálculo masivo en paralelo, vital para el entrenamiento e inferencia de redes neuronales.
    \item[Ground Truth (verdad de terreno):] conjunto de datos reales de referencia, que pueden obtenerse mediante observaciones de campo, con los que poder evaluar la calidad y precisión de las predicciones de un modelo.
    \item[Itinerario peatonal accesible (IPA):] itinerario con un flujo de paso continuo de al menos 1,80 metros de ancho y 2,20 metros de alto que permite la circulación a cualquier persona de forma segura, cómoda y autónoma (Orden TMA/851/2021, artículo 5.1).
    \item[Kanban:] metodología ágil de gestión de proyectos que utiliza un flujo de tareas por estados (pendiente, en proceso, en pruebas o finalizado), que permite visualizar y organizar el desarrollo de un trabajo o proyecto.
    \item[Movilidad reducida:] persona que tiene una limitación en su movilidad que le afecta en sus desplazamientos de manera funcional. Entre otros, incluye a usuarios/as de silla de ruedas o personas mayores.
    \item[Movilidad reducida temporal:] persona que tiene movilidad reducida durante un periodo transitorio. Se incluyen: personas con muletas, lesiones o familias con carritos de bebé.
    \item[Orden TMA/851/2021:] regulación estatal marco que establece las condiciones mínimas de accesibilidad en vías urbanas, y que define las características mínimas para considerar una vía accesible y la normativa sobre plataformas únicas. Sustituye a la Orden VIV/561/2010.
    \item[Plan de accesibilidad universal (PAU):] instrumento técnico de diagnóstico de la accesibilidad municipal, que debe incluir un informe con el listado de barreras en la vía urbana, priorización de actuaciones y un plan de actuación y presupuestario.
    \item[Producto Mínimo Viable (MVP):] mínima unidad de entrega de un producto que contiene los objetivos fundamentales y esenciales del mismo, que permite ofrecer una funcionalidad mínima en un periodo corto de tiempo.
    \item[Segmentación semántica:] técnica de visión por computador capaz de detectar clases semánticas a partir de una imagen digital (acera, calzada, vehículo, vegetación o mobiliario), clasificando píxel a píxel y utilizando máscaras de predicción.
    \item[Sistema de información geográfica (SIG/GIS):] herramienta tecnológica que gestiona información espacial y cartográfica, permitiendo integrar diferentes capas vectoriales para su análisis y representación en mapas.
    \item[Transfer learning:] técnica del deep learning que reutiliza modelos preentrenados que funcionan bien en tareas generales, que realiza un ajuste de pesos llamado \textit{fine-tuning} para la especialización de un problema específico, optimizando el tiempo de entrenamiento.
    \item[Visión computacional:] subconjunto de la inteligencia artificial especializado en el procesamiento y análisis automático de imágenes digitales, permitiendo la extracción de información semántica estructurada a partir de la naturaleza desestructurada de las imágenes.
    \item[Visor SIG (GIS):] interfaz gráfica interactiva que representa datos en un mapa, permitiendo la visualización y consulta de elementos geográficos al usuario.
\end{description}
\end{multicols}


\chapter{Bibliografía}
\begin{thebibliography}{99}

\bibitem{ayto_albacete_2021} 
Ayuntamiento de Albacete (2021) 
\textit{Guía de buenas prácticas de accesibilidad universal del Ayuntamiento de Albacete}. 
Albacete: ACCESIA y COCEMFE Albacete.

\bibitem{ayto_madrid_movilidad_2018} 
Ayuntamiento de Madrid (2018) 
«Ordenanza de Movilidad Sostenible, de 5 de octubre de 2018», 
\textit{Boletín Oficial del Ayuntamiento de Madrid}, nº 8259.

\bibitem{ayto_madrid_limpieza_2022} 
Ayuntamiento de Madrid (2022) 
«Ordenanza de Limpieza de los Espacios Públicos y Gestión de Residuos, de 26 de mayo de 2022», 
\textit{Boletín Oficial del Ayuntamiento de Madrid}, nº 9152.

\bibitem{cscae_fundacion_once_2024} 
CSCAE y Fundación ONCE (2024) 
\textit{Documenta 1.2 - Ciudad y Territorio justo: Accesibilidad universal}. 
Madrid: Consejo Superior de los Colegios de Arquitectos de España y Fundación ONCE. 
Disponible en: \url{https://observatorio2030.com/sites/default/files/2025-01/Documenta\%201.2\%20-\%20Ciudad\%20y\%20Territorio\%20justo\%20-\%20Accesibilidad\%20universal\%20\%28Informe\%20GT1.2\%29\%20\%5BAccesible\%5D_1.pdf}.

\bibitem{dacunha_2015} 
da Cunha, I. (2015) 
\textit{El trabajo de fin de grado y de máster: Redacción, defensa y publicación}. 
Barcelona: Editorial UOC.

\bibitem{fundacion_once_2022} 
Fundación ONCE (2022) 
\textit{Resumen ejecutivo: Índice de accesibilidad e inclusión. Ciudades de España 2022}. 
Madrid: Fundación ONCE. 
Disponible en: \url{https://biblioteca.fundaciononce.es/system/files/resumen_ejecutivo_indice_de_accesibilidad_e_inclusion._ciudades_de_espana_2022_def.pdf}.

\bibitem{ine_edad_2020} 
Instituto Nacional de Estadística (INE) (2022) 
\textit{Encuesta sobre Discapacidad, Autonomía Personal y Situaciones de Dependencia (EDAD 2020)}. 
Nota de prensa. Madrid: INE. 
Disponible en: \url{https://www.ine.es/prensa/edad_2020_p.pdf}.

\bibitem{rdl_1_2013} 
Jefatura del Estado (2013) 
«Real Decreto Legislativo 1/2013, de 29 de noviembre, por el que se aprueba el Texto Refundido de la Ley General de derechos de las personas con discapacidad y de su inclusión social», 
\textit{Boletín Oficial del Estado}, nº 289, págs. 95635--95673.

\bibitem{ministerio_derechos_sociales_2017} 
Ministerio de Derechos Sociales y Agenda 2030 y Fundación ONCE (2017) 
\textit{Estudio de accesibilidad universal en espacios públicos urbanizados y en la edificación en España}. 
Madrid: Observatorio de la Accesibilidad. 
Disponible en: \url{https://observatoriodelaaccesibilidad.es/wp-content/uploads/2024/04/Estudio-de-accesibilidad-universal-en-espacios-publicos-urbanizados-y-en-la-edificacion-en-espana-2017.pdf}.

\bibitem{rd_505_2007} 
Ministerio de la Presidencia (2007) 
«Real Decreto 505/2007, de 20 de abril, por el que se aprueban las condiciones básicas de accesibilidad y no discriminación de las personas con discapacidad para el acceso y utilización de los espacios públicos urbanizados y edificaciones», 
\textit{Boletín Oficial del Estado}, nº 113, págs. 20380--20387.

\bibitem{mitma_orden_tma851_2021} 
Ministerio de Transportes, Movilidad y Agenda Urbana (2021) 
«Orden TMA/851/2021, de 23 de julio, por la que se desarrollan las condiciones básicas de accesibilidad y no discriminación para el acceso y utilización de los espacios públicos urbanizados», 
\textit{Boletín Oficial del Estado}, nº 187, págs. 96328--96383.

\bibitem{oates_2006} 
Oates, B.J. (2006) 
\textit{Researching Information Systems and Computing}. 
London: Sage Publications.

\bibitem{once_pavimento_2021} 
ONCE (2021) 
\textit{Pavimento tacto-visual en el entorno urbano}. 
Madrid: Grupo Social ONCE.

\bibitem{once_plataformas_2024} 
ONCE (2024) 
\textit{Tipos de Plataformas Únicas. Condiciones de accesibilidad}. 
Madrid: Grupo Social ONCE.

\bibitem{predif_2016} 
PREDIF (2016) 
\textit{Criterios e indicadores de accesibilidad universal en entornos urbanos y turísticos}. 
Madrid: Plataforma Representativa Estatal de Personas con Discapacidad Física.

\bibitem{sala_alonso_2012} 
Sala Mozos, E. y Alonso López, F. (2012) 
\textit{La Accesibilidad Universal en los municipios: Guía para una política integral de promoción y gestión}. 
Madrid: IMSERSO - Ministerio de Sanidad, Servicios Sociales e Igualdad.

\bibitem{segittur_2021} 
SEGITTUR (2021) 
\textit{Turismo para todos: un reto para el sector y un derecho para las personas}. 
Madrid: Sociedad Mercantil Estatal para la Gestión de la Innovación y las Tecnologías Turísticas. 
Disponible en: \url{https://www.segittur.es/blog/sin-categorizar/turismo-para-todos-un-reto-para-el-sector-y-un-derecho-para-las-personas/}.

\bibitem{vazquez_ortega_2020} 
Vázquez Ortega, M. (2020) 
\textit{Accesibilidad en el 112: Diseño de un asistente virtual para personas sordas}. 
Trabajo Fin de Máster. Barcelona: Universitat Oberta de Catalunya (UOC).

\end{thebibliography}

\newpage
\appendix
\chapter{Registro de uso de inteligencia artificial generativa}

\label{anx:ia}

En este anexo se indican todos los usos de inteligencia artificial generativa realizados durante la elaboración de este trabajo, indicando la herramienta usada, el objetivo de su uso, el prompt empleado junto con los parámetros de entrada y el resultado obtenido, en caso de que sea necesario.
No se incluyen como resultados obtenidos aquellos que ya están incluidos en la memoria, sino que se hace una referencia a la sección de la memoria correspondiente.
 
\begin{fichaia}{FichaIA. Registro de trazabilidad y generación de metadatos de IA}
  \fichadato{Fecha}{30/09/2026}
  \fichadato{Herramienta}{Gemini (Google)}
  \fichadato{Fase}{WP1. Definición y objetivos}
  \fichadato{Objetivo}{Ofrecer de forma transparente la infomación sobre el uso de inteligencia artificial generativa utilizada por el autor, incluyendo los prompts y ajustes realizados.}
  \fichadato{Resultado}{Código de ficha en LaTeX (\texttt{\textbackslash begin\{fichaia\}}) normalizado que compila los comandos, parámetros y modificaciones de la sesión.}
  \fichadato{Validación}{Elaboración propia y verificación de la estructura exigida para el anexo de transparencia de la memoria del TFM.}
  \fichaprompt{A partir de la conversación o de las modificaciones que acabamos de realizar, genera la ficha de registro en formato LaTeX \textbackslash begin\{fichaia\}... (incluyendo la ficha completa del componente en Latex)}
\end{fichaia}

\begin{fichaia}{Sistema de diseño de la memoria}
  \fichadato{Fecha}{29/09/2026 (con ajustes hasta el 30/09/2026)}
  \fichadato{Herramienta}{Claude (Anthropic)~\cite{claude}}
  \fichadato{Fase}{WP1. Definición y objetivos}
  \fichadato{Objetivo}{Disponer de un sistema de diseño completo para la memoria, siguiendo los estilos base definidos por la universidad, aportando una visualización homogénea y atractiva, y facilitando al autor la generación de informes, gráficas, infografías, tablas o mapas.}
  \fichadato{Resultado}{Guía de estilo y sistema de diseño para la memoria, junto con un listado de componentes de LaTeX.}
  \fichadato{Validación}{Comprobación de la compilación con la plantilla base de la universidad; revisión visual del autor; comprobación a medida que se han ido incorporando los componentes a la memoria.}
  \fichaprompt{Estoy haciendo un trabajo fin de master del máster de ciencia de datos de la UOC, necesito un sistema de diseño para poder generar informes, gráficos, infografías, presentación de resultados, diagramas de arquitectura de redes neuronales, arquitectura y diseño de código. La temática del proyecto es la utilización de imágenes aéreas de alta resolución para extraer información de las calles mediante deep learning, segmentación semántica y visión por computador, para permitir clasificar las calles en base a su accesibilidad, para poder ayudar a generar los planes de accesibilidad con información sacda de fotos, y reducir el gasto y tiempo que lleva la inspección in situ. Te paso unas capturas del estilo definido por la universidad para la entrega de la memoria. Evidentemente el formato de la memoria es académico, pero podría llegar a tener algún cambio si así lo solicito. ¿Me puedes crear el sistema de diseño? Si tienes alguna pregunta previa, me dices}
  \fichadato{Ajustes}{Solicitud de colores accessibles; cambio a idioma castellano; solicitud de texto a dos columnas; solicitud de nuevos componentes (diagramas de flujo, conceptual y de secuencia, Kanban, estados, gráficos de barras, líneas y circular); solicitud de creación de componentes para citas y estado del arte; solicitud de componentes para comparación real frente a predicho; solicitud de componentes para notas; mejora de estilos de listas; componente para citas destacadas.}
\end{fichaia}
 
\begin{fichaia}{Figuras~\ref{fig:ipa-plano} y \ref{fig:ipa-seccion}. Itinerario peatonal accesible}
  \fichadato{Fecha}{30/09/2026}
  \fichadato{Herramienta}{Claude (Anthropic)~\cite{claude}}
  \fichadato{Fase}{WP1. Definición y objetivos, WP2.1. Análisis legal}
  \fichadato{Objetivo}{Realizar una visualización del itinerario peatonal accesible atractiva para el lector, y comprensible, para mejorar la comprensión de términos técnicos relacionados con la ley.}
  \fichadato{Resultado}{Código TikZ de las figuras~\ref{fig:ipa-plano} y \ref{fig:ipa-seccion}, en una sola figura (la división de la figura en dos fue realizada de manera manual).}
  \fichadato{Validación}{Medidas contrastadas con el texto de la Orden TMA/851/2021.}
  \fichaprompt{NEcesito que realices un mapa con medidas utilizando el componente de mapas, que defina qué es un itinerario peatona accesible según la Orden TMA/851/2021, incluyendo: ancho libre de paso (mínimo 1,80m de anchura), sin obstáculos. Altura libre de paso, mínimo 2,20m, ubicación de mobiliario en la banda exterior de la acera (con un mínimo de 40 cm de bordillo)}
  \fichaajuste{Utiliza los componentes que tienes, ¿no podrías utilizarlos?}
  \fichaajuste{ME gustaría indicar como fuente Claude, a partir de un prompt inicando qué necesito}
\end{fichaia}
 
\begin{fichaia}{Figura~\ref{fig:infografia_accesibilidad}. La accesibilidad en tres cifras}
  \fichadato{Fecha}{30/09/2026}
  \fichadato{Herramienta}{Claude (Anthropic)~\cite{claude}}
  \fichadato{Fase}{WP1. Definición y objetivos}
  \fichadato{Resultado}{Código TikZ de la infografía (figura~\ref{fig:infografia_accesibilidad}).}
  \fichadato{Validación}{Cifras aportadas por el autor a partir del INE (EDAD 2020) y del Observatorio 2030 CSCAE--Fundación ONCE (\emph{Documenta 1.2}).}
  \fichaprompt{Necesito diseñar un gráfico conceptual que muestre que en España, en el año 2020, según el Instituto Nacional de Estadística, había 4.38 millones de habitantes con algún tipo de diversidad funcional o discapacidad. Necesitaría que mostrases en una barra el número de personas y el porcentaje de personas con discapacidad reconocida (4.38 millones, el 9,5\% de la población). Y que ampliando el concepto, en otra barra, incluyendo aparte de las personas con discapacidad oficial, sumadas a las personas mayores de 65 años (un 7.5\% estimado) con problemas de movilidad y funcionales, personas con movilidad reducida temporal (un 2.5\%), y familias con carritos de bebé (un 10.5\%). Los datos de estos tres últimos son difíciles de estimar, pero representan un 30\% en total de la población (incluidas las personas con discapacidad reconocida). También quiero incluir que hay un porcentaje superior al 30\,\% en las personas mayores de 65 años (exactamente 32.3\%). Representa esto: primero en un gráfico conceptual los dos primeros valores, y como cifra de infografía el segundo valor}
  \fichaajuste{En lugar de dar el dato aproximado... me gustaría que mostrases todo en una columna para llegar al 30\% y así puedo poner la referencia a Barra 2 (\textgreater30\,\%): Población beneficiaria real de la accesibilidad universal * (Fuente: Observatorio 2030 CSCAE \& Fundación ONCE, Documenta 1.2)., en el gráfico conceptual}
  \fichaajuste{Lo mismo es mejor explorar la idea de hacer una infografía conl los tres valores}
\end{fichaia}

\begin{fichaia}{Figura~\ref{fig:proceso_actual_propuesto}. Proceso actual e inspección automatizada de accesibilidad}
  \fichadato{Fecha}{30/09/2026}
  \fichadato{Herramienta}{Gemini (Google)~\cite{gemini}}
  \fichadato{Fase}{WP1. Definición y objetivos}
  \fichadato{Objetivo}{Maquetación de la infografía del proceso manual actual frente al proceso automático propuesto, para mejorar la visualización del elemento creado por elaboración propia a partir del design system.}
  \fichadato{Validación}{Comprobación manual al generar el PDF con el resultado obtenido por Gemini, haciendo ajustes de maquetación.}
  \fichaprompt{Puedes mover un poco más abajo los términos positivos de lo propuesto y hacer que entre ACTUAL manual (in situ) bien?}
  \fichaajuste{Prefiero que no reduzcas la letra, de actual manual, mejor amplía el alto.}
  \fichaajuste{¿Me divides la infografía en 2?}
\end{fichaia}

\begin{fichaia}{Figura~\ref{fig:gantt_tfm}. Diagrama de Gantt del proyecto}
  \fichadato{Fecha}{30/09/2026}
  \fichadato{Herramienta}{Claude (Anthropic)~\cite{claude}}
  \fichadato{Fase}{WP1. Definición y objetivos}
  \fichadato{Objetivo}{Adaptar el diagrama de Gantt anteriormente elaborado de forma manual por el autor al design system de la memoria, sin modificar el contenido.}
  \fichadato{Material aportado}{Código LaTeX original del diagrama (\texttt{pgfgantt}), con las tareas, fechas e hitos definidos por el autor.}
  \fichadato{Resultado}{Código de la figura~\ref{fig:gantt_tfm}: estilo \texttt{tfmgantt}, meses en castellano, rejilla semanal, fases WP2 y WP3 agrupadas, tareas secundarias con borde discontinuo, descanso rayado y leyenda.}
  \fichadato{Validación}{Fechas, tareas e hitos comprobados frente al original; nombres de las fases agrupadas revisados por el autor; compilación con la plantilla TFUOC.}
  \fichaprompt{Transforma este cronograma a tu estilo por favor, dame el código latex para copiar y pegar}
\end{fichaia}
 
\begin{fichaia}{Tabla~\ref{tab:planificacion_tfm}. Planificación del proyecto y carga en ECTS}
  \fichadato{Fecha}{30/09/2026}
  \fichadato{Herramienta}{Claude (Anthropic)~\cite{claude}}
  \fichadato{Fase}{WP1. Definición y objetivos}
  \fichadato{Objetivo}{Adaptar la tabla de planificación anteriormente elaborada de forma manual por el autor al design system de la memoria, manteniendo su contenido.}
  \fichadato{Material aportado}{Código LaTeX original de la tabla (\texttt{tabularx}), con todos los datos definidos por el autor.}
  \fichadato{Resultado}{Código de la tabla~\ref{tab:planificacion_tfm}: carga separada en ECTS y horas con coma decimal, fechas abreviadas, descanso diferenciado, tareas secundarias marcadas con \textsuperscript{\dag}, fila de total y nota de tabla. Ajustes menores de redacción en algunas descripciones.}
  \fichadato{Validación}{Cargas, fechas y entregas comprobadas frente al original; suma total de ECTS y horas revisada por el autor; redacción de las descripciones revisada.}
  \fichaprompt{La tabla también}
\end{fichaia}

\begin{fichaia}{Bibliografía de la memoria}
  \fichadato{Fecha}{30/09/2026}
  \fichadato{Herramienta}{Claude (Anthropic)~\cite{claude}}
  \fichadato{Fase}{WP1. Definición y objetivos}
  \fichadato{Objetivo}{Transformar la bibliografía manual de formato \texttt{thebibliography} a formato BibTeX estándar, de manera más rápida y eficiente para el autor.}
  \fichadato{Material aportado}{Archivo original \texttt{09-bibliografia.tex} con entradas manuales en formato \texttt{thebibliography}.}
  \fichadato{Resultado}{Archivo \texttt{referencias.bib} con 21 entradas bibliográficas en formato BibTeX; configuración de \texttt{\textbackslash usepackage[bibliografia]\{tfm-estilo\}} y \texttt{\textbackslash addbibresource\{referencias.bib\}} en el preámbulo; sustitución del capítulo de bibliografía por \texttt{\textbackslash printbibliography[heading=bibintoc]}.}
  \fichadato{Validación}{Compilación con biber; verificación de citas en el documento; comprobación de que todas las referencias se generan correctamente en el PDF.}
  \fichaprompt{Transforma la bibliografía en referencias.bib, y carga la bibliografía en la sección en la que está}
\end{fichaia}

\end{document}
